\documentclass[12pt]{article}
\usepackage[utf8]{inputenc}
\usepackage[T1]{fontenc}
\usepackage{hyperref}
\title{Kathariya Tharu Language Resources for K-12 Curriculum Development}
\author{Compiled by Claude Sonnet 5 (claude-sonnet-5) for Andr\'as Kornai}
\date{}
\begin{document}
\maketitle

\section{Introduction}
Kathariya Tharu (ISO 639-3 tkt) is spoken in Nepal and India. The search found no dedicated Kathariya Tharu dictionary, primer or textbook. Only a sociolinguistic survey and some related-variety material turned up.

\section{Primary Educational Resources}
None were found.

\section{Linguistic and Reference Materials}
\begin{itemize}
\item ASJP language page for Tharu Kathariya. It gives about 142,000 speakers and cites Regmi (2014), \emph{A Sociolinguistic Survey of Kathariya Tharu}, Linguistic Survey of Nepal, Tribhuvan University. \url{https://asjp.clld.org/languages/THARU_KATHARIYA}
\item Webonary Tharu--Nepali dictionary. It covers Mid-Eastern Tharu, a different variety, and is useful only for comparison. \url{https://www.webonary.work/mp-tharu/overview/forward/}
\item Glosbe pages pairing ``Kathoriya Tharu'' with other languages. These are community-built and their quality is unverified. \url{https://glosbe.com/tkt/mxv}
\end{itemize}

\section{Religious and Community Resources}
None were found.

\section{Digital and Audio Resources}
None were found beyond the unverified Glosbe pages above.

\section{Curriculum Development Context}
The language is not known to be taught in schools. The Regmi (2014) survey is the first thing to obtain. SIL Nepal could be asked about word lists.
\end{document}